Evaluation will be conducted in a high-fidelity ASV simulator (e.g., VRX/Gazebo or Unity) capable of modeling hydrodynamics and detailed sensor energy consumption. We define three primary test scenarios: `Lake` (large open areas, minimal drift), `River` (narrow channels, strong directional currents), and `Coast` (dynamic waves, wind, dynamic obstacles). The system is trained on procedurally generated maps of each type and tested on held-out maps to ensure generalization.

We compare five experimental settings: (1) **Frontier Exploration**, a standard geometry-based baseline with sensors always active; (2) **Rule-Based Mission Planner**, a lawnmower coverage pattern with simple battery turn-around logic; (3) **Energy-Unaware RL**, maximizing information gain without energy penalties; (4) **Energy-Aware RL (No Adapt)**, which includes energy penalties but lacks environment type input; and (5) our **Proposed Method**, which utilizes full environment adaptation and sensor duty-cycling.

The primary metrics are **Coverage Efficiency**, defined as the area mapped ($m^2$) per Joule of energy consumed, and **Mission Success Rate**, the percentage of missions that return to base or reach a goal before battery depletion. Secondary metrics include total information gain (entropy reduction) and safety (collision rate). We also perform stress tests, such as starting missions with critical battery levels (30\%) or introducing sudden environmental shifts (e.g., transitioning from calm lake to strong river current), to evaluate robustness. Prior published results suggest that adaptive strategies significantly outperform static baselines in resource-constrained tasks (Fig.~
ef{fig:published-evidence}).
